\begin{figure*}
    \centering
    
    % 调用外层框架
    \begin{CaseStudyFrame}{Case Study: Human Intuition v.s. World Model Inference}{fig:case_study_quest}
        
        % 1. 上下文
        \ContextBar{
            \textbf{Task:} Google Quest Challenge (Multi-label Subjective QA Regression). \\
            \textbf{Data Profile:} Small scale ($N_{train} \approx 5.5k$) with 30 \textit{heavily discretized} and \textit{skewed} targets. \\
            \textbf{Data Report Insight:} Strong Question-Group structure (85\% of target variance is explained by Q-means) creates high \textbf{Leakage Risk} in random splits.
        }
        
        % 2. 左右分栏结构
        % --- 左栏：Solution 0 (Deep Learning) ---
        \begin{minipage}[t]{0.49\linewidth}
            \begin{SolutionColumn}{myblue!50!black!80}{red!10}{Solution 0: Deep Learning w/ Cross-Attention}{A}
                \small % 整体字体变大
                % Stack 去掉 texttt，使用默认字体
                \textbf{Stack:} SentenceTransformer + MultiHead-CrossAttn + AdamW \\
                \textbf{Setup:} Holdout Validation (Random Split), MSE Loss.
                \par\vspace{2pt}
                \textbf{Actual Score:} \textcolor{red!60!black}{\textbf{0.2961} \ding{116}} (Spearman $\rho$)
                
                \vspace{0.2cm} \hrule \vspace{0.2cm}
                
                % Human View (Gray)
                % Human 预测它是 Favorite (赢)，实际输了 -> 错误 (Cross)
                \HumanView{Neural Network with Cross-Attention explicitly models QA interaction. Pre-trained embeddings capture superior semantics compared to statistical features.}{Strong Favorite \textcolor{red}{\faTimes}}
                \vspace{-2pt}
                
                % WM View (Blue)
                % WM 批评它有过拟合风险 (输)，实际输了 -> 正确 (Check)
                \WMView{Critique}{Model is underconstrained for small data ($5.5k$). Random split ignores question groups, causing leakage. High overfitting risk.}{Overfitting Risk \textcolor{green!60!black}{\faCheck}}
            \end{SolutionColumn}
        \end{minipage}%
        \hfill
        % --- 右栏：Solution 1 (LightGBM) ---
        \begin{minipage}[t]{0.49\linewidth}
            \begin{SolutionColumn}{myblue!50!black!80}{green!10}{Solution 1: LightGBM Ensemble}{A}
                \small % 整体字体变大
                % Stack 去掉 texttt
                \textbf{Stack:} TF-IDF + TruncatedSVD + LightGBM \\
                \textbf{Setup:} 5-Fold Cross Validation + Weighted Ensemble Optimization.
                \par\vspace{2pt}
                \textbf{Actual Score:} \textcolor{green!40!black}{\textbf{0.3145} \ding{115}} (Spearman $\rho$)
                
                \vspace{0.2cm} \hrule \vspace{0.2cm}
                
                % Human View (Gray)
                % Human 认为它是 Weak (输)，实际赢了 -> 错误 (Cross)
                \HumanView{TF-IDF is outdated. LightGBM is typically for tabular data, not text. The method is too simple, likely to underfit complex modern natural language processing tasks.}{Weak Baseline \textcolor{red}{\faTimes}}
                \vspace{-2pt}
                
                % WM View (Blue)
                % WM 认为它是 Optimal (赢)，实际赢了 -> 正确 (Check)
                \WMView{Insight}{Robust choice. 5-Fold CV provides honest estimates. SVD captures global label families. Sample-efficient given data scarcity.}{Optimal Fit \textcolor{green!60!black}{\faCheck}}
            \end{SolutionColumn}
        \end{minipage}
        
        % 3. 底部总结
        \OutcomeFooter{Solution 1 outperformed Solution 0. The World Model correctly prioritized \textbf{Validation Rigor} and \textbf{Sample Efficiency} over Architectural Sophistication.}
        
    \end{CaseStudyFrame}
    
    \caption{\textbf{Case Study: Human Intuition vs. World Model Inference.} This example illustrates a \textbf{hidden logical conflict} where architectural sophistication (favored by humans) clashes with data constraints. By leveraging the generated \textbf{Data Report}, the World Model detects a critical mismatch between the small dataset size ($N \approx 5.5k$) and the complex neural network (Solution 0). It correctly prioritizes the robust LightGBM ensemble (Solution 1), demonstrating the ability to weigh \textit{Data-Model Fit} over pure algorithmic complexity.}

    \label{fig:case_study_quest}
\end{figure*}





\begin{figure*}[b]
    \centering
    
    % 调用外层框架
    \begin{CaseStudyFrame}{Case Study: Prioritizing Domain Fit over SOTA Architecture}{fig:case_study_tgs}
        
        % 1. 上下文
        \ContextBar{
            \textbf{Task:} TGS Salt Identification (Pixel-level Segmentation). \\
            \textbf{Data Profile:} Small scale ($101 \times 101$) grayscale seismic images with weak intensity contrast and high variability in mask coverage. \\
            \textbf{Data Report Insight:} High sensitivity to spatial artifacts; resizing introduces interpolation noise that destroys local salt boundary details. Depth data is unavailable (NaN).
        }
        
        % 2. 左右分栏结构
        % --- 左栏：Solution 0 (ViT) ---
        \begin{minipage}[t]{0.49\linewidth}
            \begin{SolutionColumn}{myblue!50!black!80}{red!10}{Solution 0: Vision Transformer}{A}
                \small
                \textbf{Stack:} ViT-B/16 (Pre-trained on ImageNet) \\
                \textbf{Setup:} Resize to $224 \times 224$, BCE + Dice Loss, attempts to use depth channel.
                \par\vspace{2pt}
                \textbf{Actual Score:} \textcolor{red!60!black}{\textbf{0.1301} \ding{116}} (mAP)
                
                \vspace{0.2cm} \hrule \vspace{0.2cm}
                
                % Human View (Gray)
                \HumanView{ViT is a top-tier SOTA model offering global context modeling. Pre-training on ImageNet guarantees strong feature extraction capabilities.}{Strong Favorite \textcolor{red}{\faTimes}}
                \vspace{-2pt}
                
                % WM View (Blue)
                \WMView{Critique}{Resizing $101 \times 101$ inputs to $224 \times 224$ risks losing fine details. Domain gap between ImageNet (RGB) and seismic data (grayscale) is severe. Incorrectly relies on missing depth data.}{Resolution Mismatch \textcolor{green!60!black}{\faCheck}}
            \end{SolutionColumn}
        \end{minipage}%
        \hfill
        % --- 右栏：Solution 1 (U-Net) ---
        \begin{minipage}[t]{0.49\linewidth}
            \begin{SolutionColumn}{myblue!50!black!80}{green!10}{Solution 1: Standard U-Net}{A}
                \small
                \textbf{Stack:} Custom U-Net with Hybrid Attention \\
                \textbf{Setup:} Native $101 \times 101$ resolution, BCE Loss, Test Time Augmentation (TTA).
                \par\vspace{2pt}
                \textbf{Actual Score:} \textcolor{green!40!black}{\textbf{0.6396} \ding{115}} (mAP)
                
                \vspace{0.2cm} \hrule \vspace{0.2cm}
                
                % Human View (Gray)
                \HumanView{Standard U-Net is an older, basic baseline. It lacks the global receptive field and advanced attention mechanisms of modern transformers.}{Weak Baseline \textcolor{red}{\faTimes}}
                \vspace{-2pt}
                
                % WM View (Blue)
                \WMView{Insight}{Preserves native spatial resolution without interpolation noise. Attention mechanisms efficiently capture multi-scale features and spatial dependencies crucial for weak-contrast segmentation.}{Optimal Fit \textcolor{green!60!black}{\faCheck}}
            \end{SolutionColumn}
        \end{minipage}
        
        % 3. 底部总结
        \OutcomeFooter{Solution 1 outperformed Solution 0. The World Model was not swayed by the ``SOTA'' status, correctly prioritizing \textbf{Native Spatial Resolution} and \textbf{Domain Fit} over superficial architectural complexity.}
        
    \end{CaseStudyFrame}
    
    \caption{\textbf{Case Study: Domain Fit vs. Architectural Sophistication.} This example highlights a ``square peg in a round hole'' scenario. While human intuition might favor the SOTA Vision Transformer (Solution 0), the World Model detects that forcing small $101 \times 101$ seismic images into ViT's required resolution introduces destructive interpolation noise. It rejects the complex model in favor of a U-Net (Solution 1) that preserves native spatial details and avoids relying on missing depth data, demonstrating skepticism towards superficially sophisticated but methodologically flawed approaches.}
    \label{fig:case_study_tgs}
\end{figure*}




\begin{figure*}
    \centering

\begin{ReportFrame}{Case Study: Verbal Data Report ($D_{rep}$) Sample for Task ``US Patent Matching''}{fig:long_report_content_1}

\#\# Data Overview

Train: 32,825 pairs; Test: 3,648 pairs. Columns: id, anchor... score (train only; in \{0.0... 1.0\}). No missing values...

106 unique 4-character CPC contexts; coverage across major sections... broad and similar in train/test.

Anchors: 733 unique; heavily reused (mean \textasciitilde 45 pairs per anchor...). Targets: 26,850 unique...

Test anchors: 100\% seen... Test targets: \textasciitilde 29\% seen... Test OOV rate \textasciitilde 12.5\%.
\newline

Why this structure matters:

The evaluation setting is anchor-centric... rewards learning anchor-specific decision boundaries and context-conditioned mappings...
\newline

\#\# Key Statistical Findings

Discrete labels concentrated at 0.25 and 0.5... Implication: strong class imbalance toward moderate similarity...

Correlation with score: char 3-5gram TF-IDF cosine \textasciitilde 0.46... Implication: surface overlap explains much variance but not all...

Phrases are very short (mean \textasciitilde 2 tokens...)... Implication: models will rely on token/subword and character patterns...

Context means vary... Implication: the mapping from lexical similarity to score is context-dependent...

Anchors average \textasciitilde 45 target pairs... Implication: each anchor induces a nontrivial local decision boundary...

Distributions are stable... Implication: generalization hinges on handling unseen targets...
\newline

\#\# Implications for Model Design

(Linking observation → modeling implication → evaluation impact)

Loss functions that ignore ordinality may mis-penalize near misses... Rank-sensitive metrics will reward monotonic mappings...

Architectures emphasizing character/subword patterns... align with dominant signal... Performance differences will emerge in tail cases...

Pairwise encoders that allow rich anchor-target interactions can learn... Approaches that exploit anchor identity will likely score well...

Models benefit from mechanisms that allow interaction between context and similarity features... Per-context performance may vary...

High-capacity sequence models may be underutilized on such short inputs... Efficiency-capacity trade-offs skew toward models effective on short spans...

Representations that degrade gracefully on unseen words... have an advantage... Robust handling of OOV will particularly improve performance...

Simple identity detection captures some easy gains... however, near-identical forms can still have scores <1.0...

Character CNNs... Token-level transformers... Independent encoders... Joint encoders... Denoising pretraining...
\newline

\#\# Summary

The dataset is anchor-centric... Surface-form similarity explains a large portion... Evaluation emphasizes generalization to new targets...

Capture strong character/subword overlap signals... Maintain robustness to unseen target tokens... Avoid over-reliance on identity heuristics... Account for label ordinality...
    
\end{ReportFrame}
    
\captionof{figure}{\textbf{Case Study: Verbal Data Report ($D_{rep}$) Sample for ``US Patent Matching''.} Generated via the \textit{Code-Execution-Verbalization} protocol, this artifact bridges the gap between raw data statistics and semantic reasoning.}
\label{fig:report_patent}

\end{figure*}




\begin{figure*}
    \centering

\begin{ReportFrame}{Case Study: Task Instruction ($I$) for Task ``Denoising Dirty Docs''}{fig:long_report_content_0}

\# Overview

\#\# Description

Optical Character Recognition (OCR) is the process of converting typed or handwritten documents into a digitized format. OCR makes previously static content editable, searchable, and easier to share.

This task focuses on improving OCR performance for degraded scanned documents. Given a dataset of noisy text images, the goal is to develop a model that removes visual noise such as stains, fading, and wrinkles, producing cleaner text images suitable for further processing and digitization.

\#\# Evaluation

Submissions are evaluated using the root mean squared error (RMSE) between the predicted cleaned pixel intensities and the ground truth grayscale pixel intensities.

\#\#\# Submission Format

Each image is represented as a list of pixel values with identifiers in the form 'image\_row\_col' (e.g. '1\_2\_1' for image 1, row 2, column 1). Intensity values range from 0 (black) to 1 (white).

\verb|```|

id,value
1\_1\_1,1
1\_2\_1,1
1\_3\_1,1
...

\verb|```|

\#\# Data

\#\#\# Dataset Description

The dataset contains two sets of images: 'train' and 'test'.  
- The 'train' set includes noisy text images and their corresponding cleaned versions ('train\_cleaned').
- The 'test' set contains only noisy images that need to be denoised.

The noise simulates real-world artifacts commonly seen in scanned documents, such as blur, stains, and faded ink. The task is to build a model that restores the test images to a clean, readable form.

\#\#\# File Description

- There are three directory corresponding to the data description above: 'train', 'train\_cleaned' and 'test'.
- The sample submission is stored in sampleSubmission.csv.
    
\end{ReportFrame}
    
\captionof{figure}{\textbf{Case Study: Task Instruction ($I$) for Task ``Denoising Dirty Docs''.} 
This example illustrates the raw natural language input $I$ as defined in Section~\ref{sec:agent_paradigm}. 
It outlines the problem context, dataset specifications, and evaluation criteria, serving as the foundational prompt that initiates the agent's solution generation process.}
\label{fig:task_instruction_denoising}

\end{figure*}